% !TEX root = ../tfg_etsiinf_NombreApellidos.tex
\chapter{Introducción} \label{chp:intro}

Este capítulo presenta el marco general del Trabajo Fin de Máster, centrado en
la aplicación de aprendizaje por refuerzo al control de vehículos aéreos no
tripulados dentro del ecosistema Aerostack2. En primer lugar se expone la
motivación del proyecto y su relevancia en el contexto actual de la robótica
aérea autónoma. A continuación se describe el contexto técnico y de
investigación en el que se enmarca el trabajo, así como el estado actual del
desarrollo realizado hasta la fecha. Posteriormente se detallan el objetivo
general y los objetivos específicos que guían la tesis. Por último, se resume
la estructura del documento.

%%%%%%%%%%%%%%%%%%%%%%%%%%%%%%%%%%%%%%%%%%%%%%%%%%%%%%%%%%%%%%%%%%%%%%%%%%%%%%%%
%%%%%%%%%%%%%%%%%%%%%%%%%%%%%%%%%%%%%%%%%%%%%%%%%%%%%%%%%%%%%%%%%%%%%%%%%%%%%%%%

\section{Motivación del proyecto} \label{sct:intro:motivacion}

Los vehículos aéreos no tripulados han pasado en pocos años de ser una
tecnología principalmente experimental a convertirse en una herramienta útil en
aplicaciones industriales, logísticas, de inspección y competición. En todos
estos escenarios, la autonomía y la capacidad de respuesta del vehículo son
aspectos críticos, especialmente cuando el dron debe seguir trayectorias
dinámicas, reaccionar ante referencias cambiantes o mantener un comportamiento
estable a altas velocidades.

Dentro de este contexto, el control clásico basado en reguladores PID sigue
siendo una solución ampliamente utilizada por su simplicidad, interpretabilidad
y robustez. Sin embargo, cuando la dinámica del sistema es fuertemente no
lineal y el régimen de operación es exigente, el ajuste manual de estos
controladores puede resultar costoso y su capacidad de generalización limitada.
Frente a ello, el aprendizaje por refuerzo ofrece una alternativa atractiva:
permitir que una política de control aprenda a actuar a partir de la interacción
con un entorno, optimizando directamente un criterio de comportamiento definido
mediante recompensas.

La motivación de este trabajo surge precisamente en la intersección entre
robótica aérea e inteligencia artificial. En el grupo CVAR
(\emph{Computer Vision and Aerial Robotics}) existe interés en desarrollar
soluciones de control cada vez más ágiles para drones, en particular en el
marco de la competición A2RL, donde el rendimiento del vehículo depende de su
capacidad para ejecutar maniobras rápidas y precisas. Para que este tipo de
aproximación sea viable, resulta imprescindible disponer de un entorno de
entrenamiento que sea fiel desde el punto de vista dinámico, pero también lo
suficientemente ligero como para permitir entrenar múltiples agentes en
paralelo.

Este Trabajo Fin de Máster responde a esa necesidad. La idea central no es
solamente entrenar un controlador con aprendizaje por refuerzo, sino construir
una infraestructura experimental que conecte de forma correcta el simulador
dinámico de Aerostack2 con una interfaz estándar de aprendizaje, de manera que
el entrenamiento, la evaluación y la comparación con soluciones convencionales
puedan realizarse de forma sistemática y reproducible.

%%%%%%%%%%%%%%%%%%%%%%%%%%%%%%%%%%%%%%%%%%%%%%%%%%%%%%%%%%%%%%%%%%%%%%%%%%%%%%%%
%%%%%%%%%%%%%%%%%%%%%%%%%%%%%%%%%%%%%%%%%%%%%%%%%%%%%%%%%%%%%%%%%%%%%%%%%%%%%%%%

\section{Contexto del proyecto} \label{sct:intro:contexto}

El trabajo se desarrolla sobre Aerostack2, un \emph{framework} para robótica
aérea basado en ROS 2 que proporciona componentes reutilizables para control,
estimación, comportamientos de vuelo e integración con diferentes plataformas y
simuladores. En particular, este proyecto utiliza el simulador dinámico de
multirrotores de Aerostack2, una herramienta sin interfaz gráfica que modela la
dinámica del vehículo y que puede visualizarse mediante RViz2. La elección de
este simulador es especialmente relevante para aprendizaje por refuerzo, ya que
su bajo coste computacional permite plantear entornos vectorizados y ejecutar
muchas instancias en paralelo, aumentando la eficiencia del muestreo durante el
entrenamiento.

La propuesta inicial del tutor establece una progresión clara. En una primera
fase es necesario comprender el funcionamiento de Aerostack2, del simulador de
multirrotores y de las interfaces ROS 2 disponibles. Sobre esa base debe
construirse un entorno de aprendizaje por refuerzo cuya entrada y salida estén
definidas de forma coherente con el problema de control planteado. Más adelante,
una vez validada la comunicación y la consistencia temporal de los datos, se
abordarán aspectos propiamente relacionados con el aprendizaje, como el diseño
de recompensas, la elección del algoritmo, el ajuste de hiperparámetros o la
comparación frente a un controlador PID convencional.

El estado actual del proyecto refleja precisamente esa primera fase de
integración. En el repositorio \texttt{rl\_uav\_aerostack2} se ha implementado un
entorno mínimo en Gymnasium, denominado \texttt{AS2TestEnv}, cuya finalidad es
verificar la conectividad entre el simulador y el entorno de aprendizaje. Este
entorno encapsula el ciclo básico \texttt{reset-step-close}, inicializa la
interfaz \texttt{DroneInterface} de Aerostack2, arma el dron, activa el modo
\emph{offboard}, ejecuta el despegue, envía comandos de velocidad lineal y lee
observaciones compuestas por posición y velocidad. Junto a ello, el repositorio
incluye scripts de lanzamiento del simulador, configuración del mundo, un
controlador PID de referencia y un script de prueba que permite comprobar el
intercambio de observaciones y acciones.

Por tanto, el trabajo realizado hasta el momento no corresponde todavía al
entrenamiento final del controlador, sino a una etapa fundacional imprescindible:
garantizar que la interfaz entre simulación, middleware robótico y entorno de
aprendizaje funciona de manera estable. Esta validación es especialmente
importante en un problema como el presente, donde la calidad del entrenamiento
depende de que el agente reciba observaciones correctas, aplique acciones en el
instante adecuado y pueda escalar posteriormente a ejecuciones paralelas.

%%%%%%%%%%%%%%%%%%%%%%%%%%%%%%%%%%%%%%%%%%%%%%%%%%%%%%%%%%%%%%%%%%%%%%%%%%%%%%%%
%%%%%%%%%%%%%%%%%%%%%%%%%%%%%%%%%%%%%%%%%%%%%%%%%%%%%%%%%%%%%%%%%%%%%%%%%%%%%%%%

\section{Objetivos} \label{sct:intro:objetivos}

El objetivo general de este Trabajo Fin de Máster es desarrollar un entorno de
aprendizaje por refuerzo sobre el simulador dinámico de multirrotores de
Aerostack2 y utilizarlo para entrenar y evaluar un controlador de alto nivel
para un dron, capaz de recibir como entrada una referencia de posición global
$[x, y, z]$ y generar como salida comandos de velocidad
$[v_x, v_y, v_z, \omega_z]$, manteniendo además una orientación en yaw alineada
con el vector velocidad.

Para alcanzar este objetivo general, el trabajo se divide en los siguientes
objetivos específicos:

\begin{itemize}

    \item[•] Familiarizarse con el ecosistema Aerostack2, su organización en ROS
    2 y el funcionamiento del simulador dinámico de multirrotores, identificando
    los componentes necesarios para el control del vehículo y el intercambio de
    información con el exterior.

    \item[•] Diseñar e implementar una primera interfaz entre el simulador y un
    entorno compatible con Gymnasium, de forma que el problema pueda abordarse
    posteriormente con herramientas estándar de aprendizaje por refuerzo.

    \item[•] Validar experimentalmente el ciclo básico de interacción del
    entorno, incluyendo inicialización, despegue, envío de acciones, recepción
    de observaciones y finalización segura de la simulación.

    \item[•] Extender esta base hacia un entorno vectorizado que permita
    ejecutar múltiples simulaciones en paralelo, aprovechando el bajo coste
    computacional del simulador para mejorar la eficiencia del entrenamiento.

    \item[•] Definir un problema de control progresivo en el que el agente
    aprenda a transformar referencias de posición en comandos de velocidad,
    incorporando explícitamente el requisito de \emph{path facing}, es decir,
    que el yaw del dron quede alineado con la dirección de movimiento.

    \item[•] Entrenar y evaluar un controlador basado en aprendizaje por
    refuerzo en diferentes escenarios de simulación con referencias
    secuenciales, y comparar su comportamiento frente a un controlador PID
    convencional utilizado como línea base.

    \item[•] Como posible extensión, estudiar una arquitectura jerárquica o en
    cascada en la que un segundo controlador aprenda a transformar velocidades
    deseadas en consignas de tasas angulares, aproximando así el problema a un
    control de menor nivel.

\end{itemize}

%%%%%%%%%%%%%%%%%%%%%%%%%%%%%%%%%%%%%%%%%%%%%%%%%%%%%%%%%%%%%%%%%%%%%%%%%%%%%%%%
%%%%%%%%%%%%%%%%%%%%%%%%%%%%%%%%%%%%%%%%%%%%%%%%%%%%%%%%%%%%%%%%%%%%%%%%%%%%%%%%

\section{Estructura del Documento} \label{sct:intro_estructura}

El documento se organiza en varios capítulos que recogen tanto el contexto del
problema como el desarrollo técnico del trabajo. Tras este capítulo
introductorio, el Capítulo 2 presenta el estado del arte y el contexto
tecnológico relacionado con el aprendizaje por refuerzo, el control de drones y
las herramientas empleadas. El Capítulo 3 describe el desarrollo realizado,
prestando especial atención a la integración entre Aerostack2, el simulador de
multirrotores y el entorno de aprendizaje implementado.

Posteriormente, el Capítulo 4 aborda el impacto del proyecto desde distintas
perspectivas relevantes para un trabajo de ingeniería. El Capítulo 5 recoge los
resultados obtenidos y la discusión asociada, incluyendo la validación del
entorno y, en fases posteriores del trabajo, la evaluación de los controladores
desarrollados. Finalmente, el documento se cierra con la bibliografía y los
anexos, donde se podrá incluir material complementario de carácter técnico.

%%%%%%%%%%%%%%%%%%%%%%%%%%%%%%%%%%%%%%%%%%%%%%%%%%%%%%%%%%%%%%%%%%%%%%%%%%%%%%%%
%%%%%%%%%%%%%%%%%%%%%%%%%%%%%%%%%%%%%%%%%%%%%%%%%%%%%%%%%%%%%%%%%%%%%%%%%%%%%%%%
